%!TEX root=../../main.tex

Now we analyze the structure of \( \calGl \) with the goal of establishing a
fully convex reformulation for deep ReLU networks.
This recursive constraint set encodes several key features of ReLU networks:
(i) weight sharing, through \( \Wll \), (ii) low-rank neurons, through the
outer-products \( \Tl_{j_1, \ldots, j_{l-1}, i_l} \otimes \Wll_{j_l} \),
(iii) the layer width, through
\( \sum_{j_l=1}^{p_l} \abs{\calI_{j_l}^{(l)}} \leq d_l \),
and (iv) activation cone membership.
Out of these constraints, (i) to (iii) are affected by the width of the
network.
We will show that these constraints vanish for sufficiently wide models.

Given two tensors
\( \Tl, \Vl \in \R^{d_1 \times p_1 \times \cdots \times p_l \times d_{l+1}} \)
we use \( \Vl \oplus \Vl \in \R^{d_1 \times p_1 \times \cdots \times 2 p_l \times d_{l+1}} \)
to denote the concatenation of the tensors across the second-to-last dimension.
Using this notation, consider the tensor constraint defined recursively by
\( \calC^{(1)} = \R^{d_1 \times d_2} \) and
\begin{equation}\label{eq:recursive-cone-constraint}
    \begin{aligned}
        \calC^{(l+1)}
         & = \bigg\{
        T^{(l+1)} \in \R^{d_1 \times p_1 \cdots \times p_{l} \times d_{l+1}}
        :
        T^{(l+1)}_{j_{l}, i_{l+1}}
        = T^{(l)}_{j_l\cdot i_{l+1}}
        - V^{(l)}_{j_l \cdot i_{l+1}}, \\
         & \hspace{5em} \text{ where }
        T^{(l)}_{j_l \cdot i_{l+1}}, V^{(l)}_{j_l \cdot i_{l+1}}
        \in \calK^{(l)}_{j_{l}}, \text{ and }
        T^{(l)} \oplus V^{(l)} \oplus \mathbf{0} \in \calC^{(l)},
        \bigg\}
    \end{aligned}
\end{equation}
where \( \mathbf{0} \) is the zero-tensor of shape
\( \rbr{d_1 \times p_1 \cdots \times p_{l-1} \times (d_{l} - p_{l} \cdot d_{l+1})} \).
As a result, the construction of \( \calCl \) requires the dimension of the
last order, \( d_{l} \), to satisfy \( d_{l} \geq 2 p_{l} d_{l+1} \);
otherwise, we could not embed \( T^{(l+1)} \) into \( T^{(l)} \).
It is straightforward to show that \( \calC^{(l)} \) is a convex cone
(\cref{lemma:convex-recursive-cone-constraint}).
More importantly, \( \calC^{(l)} \) is exactly the constraint set obtained
by making the network wide enough that the weight sharing, low-rank,
and width constraints in \( \calGl \) vanish.

\begin{restatable}[Constraint Convexification]{proposition}{constraintConvexification}%
    \label{prop:constraint-convexification}
    If \( d_l \geq 2 p_l d_{l+1} \) for every \( l \in [L]  \), then
    \( \calG^{(L)} = \calC^{(L)} \).
\end{restatable}

Although short, \cref{prop:constraint-convexification} is the critical 
proposition for our convex reformulation.
It says that if the network is sufficiently wide, the constraint set
in Problem~\ref{eq:non-convex-tensor-program} is convex and obtain a 
fully convex optimization program.

\aside{Now I characterize the minimum size solution before giving the true 
conditions for the convex reformulation to hold. But, I need to think about
this a bit more first.}
